\documentclass[a4paper,12pt]{article}
\usepackage[utf8]{inputenc}
\usepackage[T1]{fontenc}
\usepackage[ngerman]{babel}
\usepackage{amsmath,amssymb}
\usepackage{booktabs}
\usepackage{enumitem}
\usepackage[margin=2.5cm]{geometry}
\usepackage{tcolorbox}
\usepackage{listings}
\usepackage{xcolor}

\tcbuselibrary{breakable}

\newtcolorbox{merkbox}[1][]{
  colback=blue!5, colframe=blue!40,
  fonttitle=\bfseries, title={Merke}, breakable, #1
}
\newtcolorbox{analogie}[1][]{
  colback=green!5, colframe=green!40,
  fonttitle=\bfseries, title={Analogie}, breakable, #1
}

\lstset{basicstyle=\ttfamily\small, keywordstyle=\color{blue!70!black},
        commentstyle=\color{green!40!black}, stringstyle=\color{orange!80!black},
        numbers=left, numberstyle=\tiny, frame=single, breaklines=true,
        showstringspaces=false, language=Python}

\begin{document}

\begin{center}
{\LARGE\bfseries \"Ubungsblatt 4 -- Erkl\"arungen}\\[0.5em]
{\large VC Logik}
\end{center}

\bigskip

% ============================================================
% TEIL 0: Grundlagen
% ============================================================

\section*{Teil 0: Grundlagen -- alles, was wir brauchen}

Bevor wir uns an die Aufgaben setzen, schauen wir uns die Werkzeuge
an, die wir brauchen. Einige Begriffe kennen wir schon aus den
vorherigen Bl\"attern -- die frischen wir kurz auf. Andere Konzepte
sind neu und werden in Ruhe eingef\"uhrt. Am Ende von Teil~0 hast du
alle Bausteine, um die drei Aufgaben selbst zu l\"osen.

\subsection*{0.1 Aussagenlogische Formel -- kurze Auffrischung}

Eine aussagenlogische Formel ist ein Ausdruck, der aus \textbf{Variablen}
und \textbf{Operatoren} aufgebaut ist. Jede Variable steht f\"ur eine
ja/nein-Frage und kann den Wert $0$ (falsch) oder $1$ (wahr) annehmen.

\begin{center}
\renewcommand{\arraystretch}{1.4}
\begin{tabular}{c|l|l}
\toprule
\textbf{Symbol} & \textbf{Name} & \textbf{Bedeutung} \\
\midrule
$\neg A$ & Negation & "`Nicht $A$"' \\
$A \wedge B$ & Konjunktion & "`$A$ und $B$"' (beide wahr) \\
$A \vee B$ & Disjunktion & "`$A$ oder $B$"' (mind.\ eines wahr) \\
$A \to B$ & Implikation & "`Wenn $A$, dann $B$"' \\
$A \leftrightarrow B$ & Biimplikation & "`Genau dann $A$, wenn $B$"' \\
\bottomrule
\end{tabular}
\end{center}

Als Variablennamen schreibt man oft $p_1, p_2, p_3, \ldots$ (wenn man
viele braucht) oder einfach $A, B, C, \ldots$ (wenn man nur wenige
braucht).

\begin{merkbox}
Ein \textbf{Modell} einer Formel $\Phi$ ist eine Belegung der
Variablen mit $0$/$1$, die die Formel wahr macht.
Eine Formel kann $0$, $1$ oder \emph{sehr viele} Modelle haben.
\end{merkbox}

\begin{analogie}
Eine Formel ist wie eine T\"urschlie{\ss}anlage mit vielen Schaltern.
Jedes "`Modell"' ist eine Einstellung der Schalter, bei der die T\"ur
aufgeht. Manche T\"uren haben nur einen einzigen richtigen Code
(ein Modell). Andere lassen sich mit vielen verschiedenen Codes
\"offnen (viele Modelle). Wieder andere k\"onnen gar nicht ge\"offnet
werden (kein Modell -- die Formel ist unerf\"ullbar).
\end{analogie}

\subsection*{0.2 Warum Logik f\"ur "`richtige"' Probleme?}

Jetzt kommt der Punkt, der Logik richtig m\"achtig macht:

\begin{merkbox}[title={Die gro{\ss}e Idee}]
Viele Probleme aus der echten Welt lassen sich als Formel
\emph{codieren}. Die \textbf{Modelle der Formel entsprechen dann
genau den L\"osungen des Problems}. Ein Programm, das Modelle findet
(ein \emph{SAT-Solver}), l\"ost also das Problem.
\end{merkbox}

\begin{analogie}
Die Aussagenlogik ist eine Art \textbf{universelle Puzzle-Sprache}.
Egal ob Sudoku, Stundenplan, Schaltkreis-Verifikation oder
N-Damen-Problem -- alles kann man in dieselbe Sprache \"ubersetzen.
Und die Computer, die diese Sprache verstehen, sind heutzutage
unglaublich schnell (SAT-Solver l\"osen Formeln mit Millionen von
Variablen).
\end{analogie}

In diesem Blatt werden wir drei Beispiele f\"ur dieses Prinzip sehen:
Sudoku (Aufgabe~1), ein Schaltkreis (Aufgabe~2) und das
Bishop-Problem auf dem Schachbrett (Aufgabe~3).

\subsection*{0.3 Encoding-Muster f\"ur Kombinatorik}

Wie codiert man ein Problem? Es gibt ein wiederkehrendes Rezept:

\begin{merkbox}[title={Strategie: Problem $\to$ Formel}]
\begin{enumerate}[nosep]
\item \textbf{Eine Bool'sche Variable pro Entscheidung.} Beispiele:
      "`Steht Ziffer $n$ in Feld $(i,j)$?"', "`Steht ein L\"aufer auf
      Feld $(i,j)$?"', "`Hat Draht $o_1$ den Wert $1$?"'.
\item \textbf{Jede Regel des Problems als Teilformel.} "`Mindestens
      einer muss"' wird zu einer $\vee$-Verkn\"upfung. "`H\"ochstens
      einer darf"' wird zu einer UND-Verkn\"upfung von
      Ausschluss-Paaren.
\item \textbf{Alle Teilformeln mit $\wedge$ verbinden.} Die
      Gesamtformel ist die Konjunktion aller Regeln -- alle Regeln
      m\"ussen gleichzeitig gelten.
\end{enumerate}
\end{merkbox}

\begin{analogie}
Wir \"ubersetzen das Problem wie einen \textbf{Anforderungskatalog}.
Ein Bauherr h\"atte ihn so geschrieben:
\begin{itemize}[nosep]
\item "`In jedem Raum muss mindestens eine Lampe sein."'
\item "`In jedem Raum darf h\"ochstens eine Hauptlampe sein."'
\item "`Jeder Schalter geh\"ort zu genau einer Lampe."'
\end{itemize}
Das ist schon fast Logik: "`mindestens eine"', "`h\"ochstens eine"',
"`genau eine"'. Diese drei Muster kehren immer wieder.
\end{analogie}

\subsection*{0.4 At-most-one (AMO) -- das wichtigste Werkzeug}

Dieses Muster brauchen wir in allen drei Aufgaben. Es ist so wichtig,
dass es einen eigenen Namen hat.

\begin{merkbox}
\textbf{At-most-one (AMO).} Von $k$ Variablen $x_1, x_2, \ldots, x_k$
soll \emph{h\"ochstens eine} wahr sein. Als Formel mit
paarweiser Enkodierung:
\[
\text{AMO}(x_1, \ldots, x_k)
\;\equiv\;
\bigwedge_{1 \le i < j \le k} \neg(x_i \wedge x_j).
\]
Anzahl Klauseln: $\binom{k}{2} = \frac{k(k-1)}{2}$.
\end{merkbox}

Warum funktioniert das? Die AMO-Eigenschaft "`h\"ochstens eine wahr"'
ist gleichbedeutend mit: "`es gibt kein Paar, in dem beide wahr
sind"'. Die Teilformel $\neg(x_i \wedge x_j)$ verbietet genau dieses
Paar. Wenn wir \emph{alle} Paare verbieten, kann nirgends ein Paar
auftauchen -- also sind entweder $0$ oder $1$ Variablen wahr.

\textbf{Beispiel mit $k = 3$.} Bei drei Variablen $x_1, x_2, x_3$
brauchen wir $\binom{3}{2} = 3$ Klauseln:
\[
\neg(x_1 \wedge x_2) \;\wedge\; \neg(x_1 \wedge x_3) \;\wedge\;
\neg(x_2 \wedge x_3).
\]
Probier es aus: Wenn $x_1 = x_2 = 1$, dann ist die erste Klausel
falsch, also die ganze Formel falsch. Gut, verboten. Wenn nur
$x_1 = 1$ ist (und die anderen $0$), dann ist jede Klausel der Form
$\neg(1 \wedge 0) = \neg 0 = 1$, also wahr. Passt, erlaubt.

\begin{analogie}
Aus dem Alltag: "`Von den drei Kuchenst\"ucken darf h\"ochstens eines
gegessen werden."' Das hei{\ss}t: nicht Kuchen-A \emph{und} Kuchen-B
zugleich, nicht Kuchen-A \emph{und} Kuchen-C zugleich, nicht Kuchen-B
\emph{und} Kuchen-C zugleich. Drei Ausschl\"usse, fertig. Das ist die
AMO-Enkodierung.
\end{analogie}

\subsection*{0.5 At-least-one und Exactly-one}

Zwei Geschwister von AMO:

\begin{merkbox}
\textbf{At-least-one (ALO).} "`Mindestens eine ist wahr"':
\[
\text{ALO}(x_1, \ldots, x_k) \;\equiv\; x_1 \vee x_2 \vee \ldots \vee x_k.
\]
Das ist eine einzige gro{\ss}e Klausel.
\medskip

\textbf{Exactly-one (EO).} "`Genau eine ist wahr"':
\[
\text{EO}(x_1, \ldots, x_k) \;\equiv\; \text{ALO}(\ldots) \;\wedge\;
\text{AMO}(\ldots).
\]
Also "`mindestens eine"' \emph{und} "`h\"ochstens eine"' zusammen.
\end{merkbox}

\begin{analogie}
ALO ist: "`Mindestens ein Fenster im Haus ist offen."' Es reicht
schon, wenn \emph{irgendeins} offen ist.
\medskip

EO ist: "`Es ist genau eine Tasche am Haken."' Nicht null, nicht
zwei, genau eine. Also braucht man zwei Einzelregeln gleichzeitig:
mindestens eine (ALO) und h\"ochstens eine (AMO).
\end{analogie}

\subsection*{0.6 Schaltkreise als Formel}

F\"ur Aufgabe~2: Ein Schaltkreis besteht aus \textbf{Gates}
(Bauteilen), die durch Dr\"ahte verbunden sind. Jedes Gate hat
Eing\"ange und einen Ausgang. Der Ausgang h\"angt durch eine fest
vorgegebene Regel von den Eing\"angen ab.

\begin{merkbox}[title={Strategie: Schaltkreis $\to$ Formel}]
F\"ur \emph{jeden} Draht (insbesondere den Ausgang jedes Gates)
schreibt man eine \"Aquivalenz
\[
\text{(Draht-Variable)} \;\leftrightarrow\; \text{(Funktion der Eingangs-Variablen)}.
\]
Die Gesamtformel ist die Konjunktion all dieser \"Aquivalenzen.
\end{merkbox}

\begin{analogie}
Jedes Gate ist wie ein kleiner Mitarbeiter, der eine feste Regel
befolgt. Ein UND-Gate sagt: "`Wenn beide Inputs 1 sind, gebe ich
1 aus, sonst 0."' Um den ganzen Betrieb zu verstehen, sammeln wir
die Regel jedes einzelnen Mitarbeiters (jede Gate-Regel) und
kleben sie alle mit UND zusammen. Nur wenn \emph{jeder} Mitarbeiter
seine Regel einh\"alt, ist der Betrieb korrekt.
\end{analogie}

\textbf{Wichtig:} Es gibt sichtbare Variablen (Eingang/Ausgang des
ganzen Schaltkreises) und \textbf{Hilfsvariablen} f\"ur interne
Dr\"ahte. Die Hilfsvariablen bekommen zus\"atzliche Namen wie $o_1,
o_2, \ldots$, damit man jede einzelne Gate-Regel sauber hinschreiben
kann.

\subsection*{0.7 Erf\"ullbarkeitstester -- der Fragen-Sklave}

Aus dem letzten Blatt: Ein \textbf{Erf\"ullbarkeitstester} (SAT-Solver)
ist ein Programm, das zu einer gegebenen Formel $\Phi$ entscheidet,
ob sie ein Modell hat (SAT, erf\"ullbar) oder nicht (UNSAT,
unerf\"ullbar).

\begin{merkbox}[title={Zwei wichtige Fragen-Muster}]
\begin{itemize}[nosep]
\item \textbf{"`Kann $X$ gelten?"'} $\Rightarrow$ Pr\"ufe
      $\Phi \wedge X$ auf SAT.
      Erf\"ullbar $\Rightarrow$ Ja, $X$ ist m\"oglich.
      Unerf\"ullbar $\Rightarrow$ Nein, unter $\Phi$ ist $X$
      ausgeschlossen.
\item \textbf{"`Gilt immer $Y$?"'} $\Rightarrow$ Pr\"ufe
      $\Phi \wedge \neg Y$ auf SAT.
      Unerf\"ullbar $\Rightarrow$ Ja, $Y$ gilt immer.
      Erf\"ullbar $\Rightarrow$ Nein, Gegenbeispiel gefunden.
\end{itemize}
\end{merkbox}

\begin{analogie}
Ein Erf\"ullbarkeitstester ist wie ein treuer Diener, der nur eine
einzige Frage versteht: "`Gibt es f\"ur diese Formel eine Belegung,
die sie wahr macht?"' Alle anderen Fragen \"uber Logik-Formeln
lassen sich auf diese eine Frage reduzieren -- man muss die Frage
nur geschickt umformulieren, wie oben beschrieben.
\end{analogie}

\textbf{Warum funktioniert "`Gilt immer $Y$?"' \"uber $\neg Y$?}
"`$Y$ gilt immer (unter $\Phi$)"' bedeutet: In jedem Modell von
$\Phi$ ist $Y$ wahr. Das Gegenteil w\"are: Es gibt ein Modell von
$\Phi$, in dem $Y$ falsch ist -- also ein Modell von $\Phi \wedge \neg Y$.
Findet der Tester \emph{kein} solches Modell (UNSAT), dann gilt
$Y$ tats\"achlich immer.

\subsection*{0.8 Schach und Diagonalen}

F\"ur Aufgabe~3: Ein Schachbrett hat $n$ Zeilen und $n$ Spalten.
Die Felder bezeichnen wir mit $(i, j)$, wobei $i$ die Zeilennummer
und $j$ die Spaltennummer ist, beide von $1$ bis $n$.

\begin{merkbox}[title={Schachfiguren -- Bewegungsmuster}]
\begin{itemize}[nosep]
\item \textbf{Rook (Turm):} bewegt sich entlang Zeilen und Spalten
      (horizontal/vertikal).
\item \textbf{Bishop (L\"aufer):} bewegt sich entlang Diagonalen
      (schr\"ag).
\item \textbf{Queen (Dame):} bewegt sich wie Rook UND Bishop
      kombiniert.
\end{itemize}
\end{merkbox}

F\"ur den L\"aufer (Bishop) brauchen wir eine klare Definition von
"`Diagonale"':

\begin{merkbox}[title={Diagonalen}]
\begin{itemize}[nosep]
\item \textbf{Hauptdiagonale:} Alle Felder $(i,j)$ mit demselben Wert
      $i - j$. Die f\"uhrt von links oben nach rechts unten.
      Beispiel: $(1,1), (2,2), (3,3), (4,4)$ haben alle $i-j = 0$.
\item \textbf{Gegendiagonale:} Alle Felder $(i,j)$ mit demselben Wert
      $i + j$. Die f\"uhrt von rechts oben nach links unten.
      Beispiel: $(1,4), (2,3), (3,2), (4,1)$ haben alle $i+j = 5$.
\end{itemize}
Auf einem $n \times n$-Brett gibt es jeweils $2n - 1$ Hauptdiagonalen
und $2n - 1$ Gegendiagonalen.
\end{merkbox}

\begin{analogie}
Diagonalen sind wie \textbf{Stra{\ss}en, die quer \"uber das Brett
f\"uhren}. Jedes Feld liegt an einer Kreuzung von genau zwei
Stra{\ss}en: einer Hauptdiagonale und einer Gegendiagonale. Der
L\"aufer "`bewohnt"' also zwei Stra{\ss}en gleichzeitig. Zwei L\"aufer
k\"onnen sich nur dann schlagen, wenn sie auf derselben Stra{\ss}e
stehen -- also dieselbe Hauptdiagonale (gleiches $i-j$) oder
dieselbe Gegendiagonale (gleiches $i+j$) teilen.
\end{analogie}

\textbf{Klein-Beispiel $n = 4$.} Die Werte von $i - j$ liegen
zwischen $-(n-1) = -3$ und $n-1 = 3$, also
$i - j \in \{-3, -2, -1, 0, 1, 2, 3\}$ -- das sind $7 = 2n-1$
Hauptdiagonalen. Analog $i + j \in \{2, 3, 4, 5, 6, 7, 8\}$ -- auch
$7$ Gegendiagonalen. Schauen wir uns die Hauptdiagonale $i - j = 0$
an: Das sind die Felder $(1,1), (2,2), (3,3), (4,4)$, die vom linken
oberen Eck bis zum rechten unteren Eck verlaufen.

\newpage

% ============================================================
% AUFGABE 1
% ============================================================

\section{Aufgabe 1: Sudoku als Formel}

\subsection*{Aufgabe in einfachen Worten}

Wir haben ein $4 \times 4$-Sudoku-Feld. Jede Zelle soll eine Ziffer
von $1$ bis $4$ enthalten. Zusatzregeln: In jeder Zeile, in jeder
Spalte und in jeder der vier $2 \times 2$-Regionen darf jede Ziffer
\emph{h\"ochstens einmal} vorkommen. Gesucht: eine aussagenlogische
Formel $\Phi$, deren Modelle genau den g\"ultigen Sudoku-L\"osungen
entsprechen.

\begin{analogie}
Denk an Sudoku wie an ein \textbf{kleines R\"atsel}, bei dem du mit
vier Stempeln (mit den Ziffern 1, 2, 3, 4) auf 16 Felder stempeln
sollst. Die Regeln: In jeder Zeile, Spalte und Region darfst du jeden
Stempel h\"ochstens einmal benutzen. Unsere Aufgabe ist es, diese
Regeln so in Logik zu gie{\ss}en, dass jede g\"ultige Stempel-Anordnung
einer wahren Belegung unserer Formel entspricht.
\end{analogie}

\subsection*{Schritt 1: Die Variablen einf\"uhren}

Pro Feld gibt es $4$ M\"ogliche Ziffern. Und pro Ziffer gibt es $16$
Felder, auf denen sie stehen k\"onnte. Wir brauchen eine
Bool'sche Variable f\"ur jede \emph{Kombination}:

\begin{merkbox}[title={Die Variablen}]
F\"ur $i, j, n \in \{1, 2, 3, 4\}$ definieren wir
\[
x_{i,j,n} \;:=\; \text{``Feld $(i,j)$ enth\"alt die Ziffer $n$''}.
\]
Das ergibt $4 \cdot 4 \cdot 4 = 64$ Variablen.
\end{merkbox}

\begin{analogie}
Genauso wie in Blatt~2 die Lichtschalter-Analogie: Jede Variable ist
ein kleiner Schalter, der "`an"' (1) oder "`aus"' (0) steht. Der
Schalter $x_{1,1,3}$ steht auf "`an"', wenn im Feld in der oberen
linken Ecke die Ziffer $3$ geschrieben wird. Steht der Schalter auf
"`aus"', dann steht dort irgendetwas anderes.
\end{analogie}

\textbf{Kleines Beispiel.} Wenn das Sudoku so aussieht:
\[
\begin{array}{|c|c|c|c|}
\hline
1 & 2 & 3 & 4 \\ \hline
3 & 4 & 1 & 2 \\ \hline
2 & 1 & 4 & 3 \\ \hline
4 & 3 & 2 & 1 \\ \hline
\end{array}
\]
Dann ist $x_{1,1,1} = 1$ (im Feld oben links steht die 1), $x_{1,1,2}
= 0, x_{1,1,3} = 0, x_{1,1,4} = 0$ (dort steht keine 2, 3 oder 4).
Analog $x_{1,2,2} = 1$ (im zweiten Feld der ersten Zeile steht die 2).

\subsection*{Schritt 2: Feld-Layout und Regionen}

Das $4 \times 4$-Brett wird in vier $2 \times 2$-Regionen geteilt
(wie bei "`gro{\ss}em"' $9 \times 9$-Sudoku mit $3 \times 3$-Regionen).

\begin{center}
\renewcommand{\arraystretch}{1.5}
\begin{tabular}{|cc|cc|}
\hline
$(1,1)$ & $(1,2)$ & $(1,3)$ & $(1,4)$ \\
$(2,1)$ & $(2,2)$ & $(2,3)$ & $(2,4)$ \\ \hline
$(3,1)$ & $(3,2)$ & $(3,3)$ & $(3,4)$ \\
$(4,1)$ & $(4,2)$ & $(4,3)$ & $(4,4)$ \\ \hline
\end{tabular}
\end{center}

Die vier Regionen:
\begin{itemize}[nosep]
\item $R_1 = \{(1,1), (1,2), (2,1), (2,2)\}$ \quad (links oben)
\item $R_2 = \{(1,3), (1,4), (2,3), (2,4)\}$ \quad (rechts oben)
\item $R_3 = \{(3,1), (3,2), (4,1), (4,2)\}$ \quad (links unten)
\item $R_4 = \{(3,3), (3,4), (4,3), (4,4)\}$ \quad (rechts unten)
\end{itemize}

\subsection*{Schritt 3: Constraints schrittweise aufbauen}

Jetzt \"ubersetzen wir die Sudoku-Regeln nacheinander in Logik. Wir
stellen uns vor, dass wir einen Anforderungskatalog abarbeiten --
eine Regel nach der anderen.

\begin{analogie}
Wie beim \textbf{Kuchenbacken nach Rezept}: Erst Mehl, dann Zucker,
dann Eier. Jede einzelne Zutat ist eine kleine Formel. Am Ende
r\"uhren wir alles zusammen (mit $\wedge$) und erhalten den fertigen
Teig (die Gesamtformel).
\end{analogie}

\subsubsection*{(C1) Jedes Feld enth\"alt mindestens eine Ziffer}

Das ist ein \textbf{At-least-one}-Constraint: Die Variablen f\"ur ein
und dasselbe Feld ($x_{i,j,1}, x_{i,j,2}, x_{i,j,3}, x_{i,j,4}$)
sollen nicht alle $0$ sein.

\[
C_1 \;\equiv\;
\bigwedge_{i=1}^{4} \bigwedge_{j=1}^{4}
\bigl(x_{i,j,1} \vee x_{i,j,2} \vee x_{i,j,3} \vee x_{i,j,4}\bigr).
\]

F\"ur jedes der $16$ Felder eine Disjunktion mit $4$ Literalen.
Insgesamt also $16$ solche Disjunktionen, die alle mit $\wedge$
verkn\"upft sind.

\textbf{Konkretes Beispiel f\"ur Feld $(1,1)$:}
\[
x_{1,1,1} \vee x_{1,1,2} \vee x_{1,1,3} \vee x_{1,1,4}.
\]
Das hei{\ss}t: "`Im Feld $(1,1)$ steht mindestens eine der Ziffern
$1$, $2$, $3$, $4$."' F\"ur jedes andere Feld analog.

\subsubsection*{(C2) Jedes Feld enth\"alt h\"ochstens eine Ziffer}

Hier kommt \textbf{AMO} zum Einsatz: F\"ur ein Feld $(i,j)$ soll nicht
gleichzeitig "`hier steht Ziffer $n$"' und "`hier steht Ziffer $m$"'
gelten, wenn $n \ne m$.

\[
C_2 \;\equiv\;
\bigwedge_{i=1}^{4} \bigwedge_{j=1}^{4}
\bigwedge_{1 \le n < m \le 4}
\neg\bigl(x_{i,j,n} \wedge x_{i,j,m}\bigr).
\]

Pro Feld gibt es $\binom{4}{2} = 6$ Paare $(n,m)$, also $6$ Klauseln.
Insgesamt $16 \cdot 6 = 96$ Klauseln.

\textbf{Konkretes Beispiel f\"ur Feld $(1,1)$:} Die $6$ Paare sind
$(1,2), (1,3), (1,4), (2,3), (2,4), (3,4)$.
\begin{align*}
& \neg(x_{1,1,1} \wedge x_{1,1,2}) \;\wedge\; \neg(x_{1,1,1} \wedge x_{1,1,3}) \\
& \wedge\; \neg(x_{1,1,1} \wedge x_{1,1,4}) \;\wedge\; \neg(x_{1,1,2} \wedge x_{1,1,3}) \\
& \wedge\; \neg(x_{1,1,2} \wedge x_{1,1,4}) \;\wedge\; \neg(x_{1,1,3} \wedge x_{1,1,4}).
\end{align*}

\begin{merkbox}
(C1) und (C2) zusammen erzwingen: Jedes Feld enth\"alt \textbf{genau}
eine Ziffer. Das ist der Exactly-one-Constraint -- aber wir haben ihn
durch Kombination von ALO und AMO realisiert.
\end{merkbox}

\subsubsection*{(C3) Jede Ziffer h\"ochstens einmal pro Zeile}

Jetzt wechselt die Perspektive: Wir fixieren eine Zeile $i$ und eine
Ziffer $n$ und schauen auf die $4$ Variablen $x_{i,1,n}, x_{i,2,n},
x_{i,3,n}, x_{i,4,n}$. Davon darf h\"ochstens eine $1$ sein -- wieder
AMO.

\[
C_3 \;\equiv\;
\bigwedge_{i=1}^{4} \bigwedge_{n=1}^{4}
\bigwedge_{1 \le j < k \le 4}
\neg\bigl(x_{i,j,n} \wedge x_{i,k,n}\bigr).
\]

Pro (Zeile, Ziffer)-Kombi gibt es $\binom{4}{2} = 6$ Paare. Insgesamt
$4 \cdot 4 \cdot 6 = 96$ Klauseln.

\textbf{Konkretes Beispiel f\"ur Zeile $1$ und Ziffer $3$:}
\begin{align*}
& \neg(x_{1,1,3} \wedge x_{1,2,3}) \;\wedge\; \neg(x_{1,1,3} \wedge x_{1,3,3}) \\
& \wedge\; \neg(x_{1,1,3} \wedge x_{1,4,3}) \;\wedge\; \neg(x_{1,2,3} \wedge x_{1,3,3}) \\
& \wedge\; \neg(x_{1,2,3} \wedge x_{1,4,3}) \;\wedge\; \neg(x_{1,3,3} \wedge x_{1,4,3}).
\end{align*}

Diese Klauseln sagen: "`In Zeile $1$ gibt es nicht zwei Felder mit
der Ziffer $3$."'

\subsubsection*{(C4) Jede Ziffer h\"ochstens einmal pro Spalte}

Analog zu (C3), aber mit fixer Spalte $j$ und variierender Zeile.

\[
C_4 \;\equiv\;
\bigwedge_{j=1}^{4} \bigwedge_{n=1}^{4}
\bigwedge_{1 \le i < k \le 4}
\neg\bigl(x_{i,j,n} \wedge x_{k,j,n}\bigr).
\]

Auch hier $4 \cdot 4 \cdot 6 = 96$ Klauseln.

\subsubsection*{(C5) Jede Ziffer h\"ochstens einmal pro Region}

Und zum Schluss das gleiche Muster f\"ur die vier Regionen:

\[
C_5 \;\equiv\;
\bigwedge_{R \in \{R_1, R_2, R_3, R_4\}}
\bigwedge_{n=1}^{4}
\bigwedge_{\substack{(i,j), (k,l) \in R \\ (i,j) < (k,l)}}
\neg\bigl(x_{i,j,n} \wedge x_{k,l,n}\bigr).
\]

(Das Paar $(i,j) < (k,l)$ meint: wir nehmen jedes ungeordnete Paar nur
einmal, z.\,B.\ in lexikografischer Ordnung.) Pro (Region,
Ziffer)-Kombi gibt es $\binom{4}{2} = 6$ Paare, insgesamt wieder
$4 \cdot 4 \cdot 6 = 96$.

\textbf{Konkretes Beispiel f\"ur Region $R_1 = \{(1,1), (1,2), (2,1),
(2,2)\}$ und Ziffer $2$:}
\begin{align*}
& \neg(x_{1,1,2} \wedge x_{1,2,2}) \;\wedge\; \neg(x_{1,1,2} \wedge x_{2,1,2}) \\
& \wedge\; \neg(x_{1,1,2} \wedge x_{2,2,2}) \;\wedge\; \neg(x_{1,2,2} \wedge x_{2,1,2}) \\
& \wedge\; \neg(x_{1,2,2} \wedge x_{2,2,2}) \;\wedge\; \neg(x_{2,1,2} \wedge x_{2,2,2}).
\end{align*}

\subsection*{Schritt 4: Die Gesamtformel}

Alles zusammen:
\[
\boxed{\;\Phi \;\equiv\; C_1 \wedge C_2 \wedge C_3 \wedge C_4 \wedge C_5\;}
\]

\begin{center}
\renewcommand{\arraystretch}{1.3}
\begin{tabular}{l|l|r}
\toprule
\textbf{Constraint} & \textbf{Bedeutung} & \textbf{Klauseln} \\
\midrule
$C_1$ & Mindestens eine Ziffer pro Feld & $16$ \\
$C_2$ & H\"ochstens eine Ziffer pro Feld (AMO) & $96$ \\
$C_3$ & Keine Ziffer zweimal pro Zeile (AMO) & $96$ \\
$C_4$ & Keine Ziffer zweimal pro Spalte (AMO) & $96$ \\
$C_5$ & Keine Ziffer zweimal pro Region (AMO) & $96$ \\
\midrule
\textbf{Gesamt} & & $\mathbf{400}$ \\
\bottomrule
\end{tabular}
\end{center}

Die Formel $\Phi$ hat also $400$ Klauseln \"uber $64$ Variablen. Jedes
Modell dieser Formel entspricht einer g\"ultigen Sudoku-L\"osung.

\begin{merkbox}[title={Warum "`h\"ochstens"' statt "`genau"'?}]
Die Aufgabenstellung sagt bewusst nur "`h\"ochstens"', nicht
"`genau"'. Dadurch erlaubt die Formel auch unvollst\"andig gef\"ullte
Bretter -- zum Beispiel das leere Brett (keine Ziffern \"uberall) ist
kein Modell, weil (C1) fordert, dass jedes Feld gef\"ullt ist. Aber
wenn man (C1) weglie{\ss}e, w\"are auch das leere Brett ein Modell.
Mit (C1)+(C2) ist jedes Feld trotzdem mit genau einer Ziffer belegt,
und die Einschr\"ankungen (C3)-(C5) sorgen daf\"ur, dass keine Zeile,
Spalte oder Region doppelte Ziffern hat.
\end{merkbox}

\begin{merkbox}[title={Goldene Regel des Encodings}]
Dieses Muster funktioniert f\"ur fast jedes kombinatorische Problem:
\begin{itemize}[nosep]
\item N-Damen-Problem: $n^2$ Variablen f\"ur "`Dame auf Feld $(i,j)$"',
      AMO pro Zeile/Spalte/Diagonale.
\item Graph-F\"arbung: Variable $c_{v,n}$ f\"ur "`Knoten $v$ hat Farbe
      $n$"', EO pro Knoten, Ausschluss bei Nachbarknoten.
\item Aufzugs-Scheduling: Variable pro (Person, Stockwerk, Zeit), EO
      pro Person zu jedem Zeitpunkt, Kapazit\"ats-Constraints pro
      Aufzug.
\end{itemize}
Immer das gleiche Rezept: \textbf{Variablen pro Entscheidung +
Constraints pro Gruppe}.
\end{merkbox}

\newpage

% ============================================================
% AUFGABE 2
% ============================================================

\section{Aufgabe 2: Schaltkreis als Formel}

\subsection*{Aufgabe in einfachen Worten}

Wir haben einen Schaltkreis mit drei Eing\"angen $A, B, C$ und drei
Ausg\"angen $D, E, F$. Dazwischen gibt es "`Gates"' (Bauteile), die
aus ihren Eingangssignalen einen Ausgang berechnen. Der Clou bei
dieser Aufgabe: Die Gates rechnen nicht mit klassischen UND/ODER,
sondern mit \emph{Summen-Bedingungen} -- zum Beispiel "`Output $= 1$
genau dann, wenn die Summe der Inputs gr\"o{\ss}er als $1$ ist"'.

Gesucht: eine aussagenlogische Formel, deren Modelle genau den
zul\"assigen Belegungen $(A, B, C, D, E, F)$ entsprechen. Au{\ss}erdem
sollen wir drei Fragen zum Schaltkreis beantworten, indem wir die
Formel einem Erf\"ullbarkeitstester vorlegen.

\subsection*{Schritt 1: Gate-Semantik verstehen}

Die Notation $[> k]$ oder $[= k]$ auf einem Gate bedeutet: Der Output
ist $1$ genau dann, wenn die \textbf{Summe der Inputs}
(als Zahlen gelesen, also bin\"ar $0/1$) diese Bedingung erf\"ullt.

\begin{merkbox}[title={Summen-Bedingungen als Bool'sche Ausdr\"ucke}]
F\"ur bin\"are Inputs (je $0$ oder $1$) gilt:
\begin{itemize}[nosep]
\item $[>1](A, B)$: $A + B > 1$ $\Leftrightarrow$ beide sind $1$
      $\Leftrightarrow$ $A \wedge B$.
\item $[=1](A, C)$: $A + C = 1$ $\Leftrightarrow$ genau einer ist $1$
      $\Leftrightarrow$ $(A \wedge \neg C) \vee (\neg A \wedge C)$.
      Das ist \emph{XOR}.
\item $[>0](B, C)$: $B + C > 0$ $\Leftrightarrow$ mindestens einer ist
      $1$ $\Leftrightarrow$ $B \vee C$.
\item $[=0](x)$: $x = 0$ $\Leftrightarrow$ $\neg x$.
\item $[>2](o_1, o_2, o_3)$: Summe von drei bin\"aren Werten gr\"o{\ss}er
      als $2$ $\Leftrightarrow$ alle drei sind $1$ $\Leftrightarrow$
      $o_1 \wedge o_2 \wedge o_3$.
\item $[>1](o_2, o_3)$: $o_2 + o_3 > 1$ $\Leftrightarrow$ beide sind
      $1$ $\Leftrightarrow$ $o_2 \wedge o_3$.
\end{itemize}
\end{merkbox}

\begin{analogie}
Ein Summen-Gate ist wie eine \textbf{Abstimmung}. $[>1]$ bedeutet
"`mehr als ein Ja-Stimmer"', $[=0]$ bedeutet "`kein Ja-Stimmer"',
$[>2]$ bei drei Stimmern bedeutet "`alle stimmen ja"' usw. Der Output
ist $1$, wenn das Abstimmungsergebnis die Vorgabe erf\"ullt.
\end{analogie}

\textbf{Warum ist bei zwei bin\"aren Inputs $[>1]$ dasselbe wie
$\wedge$?} Die m\"oglichen Summen sind $0$, $1$ oder $2$. "`$> 1$"'
hei{\ss}t also $= 2$, also beide $= 1$. Das ist die UND-Verkn\"upfung.

\textbf{Warum ist $[=1]$ bei zwei Inputs XOR?} Die Summe ist $0$
(beide $0$), $1$ (genau einer $1$) oder $2$ (beide $1$). "`$=1$"'
trifft nur den mittleren Fall zu. Das ist per Definition das
Exclusive-Or.

\subsection*{Schritt 2: Verkabelung ablesen}

Schauen wir uns den Schaltkreis genau an. Der Input $A$ geht in das
Gate oben links, zugleich in das mittlere Gate. $B$ geht in das
Gate oben links und in das untere Gate. $C$ geht in das mittlere und
in das untere Gate.

Die drei linken Gates produzieren Hilfsoutputs, die wir $o_1, o_2,
o_3$ nennen. Diese Outputs gabeln sich und gehen zum Gate "`oben
rechts"' (berechnet $D$), zum Gate "`in der Mitte"' (berechnet $E$)
und zum Gate "`unten rechts"' (berechnet $F$).

\begin{center}
\renewcommand{\arraystretch}{1.3}
\begin{tabular}{l|l|l|l}
\toprule
\textbf{Gate} & \textbf{Inputs} & \textbf{Bedingung} & \textbf{Output} \\
\midrule
$[>1]$ (oben links) & $A, B$ & $A + B > 1$ & $o_1$ \\
$[=1]$ (mitte links) & $A, C$ & $A + C = 1$ & $o_2$ \\
$[>0]$ (unten links) & $B, C$ & $B + C > 0$ & $o_3$ \\
$[=0]$ (oben rechts) & $o_1$ & $o_1 = 0$ & $D$ \\
$[>2]$ (mitte rechts) & $o_1, o_2, o_3$ & Summe $> 2$ & $E$ \\
$[>1]$ (unten rechts) & $o_2, o_3$ & $o_2 + o_3 > 1$ & $F$ \\
\bottomrule
\end{tabular}
\end{center}

\begin{merkbox}
Die drei Hilfsvariablen $o_1, o_2, o_3$ sind interne Dr\"ahte -- sie
treten an den "`Au{\ss}enklemmen"' des Schaltkreises nicht auf, aber
wir brauchen sie, um die Formel sauber aufzuschreiben. Die
"`sichtbaren"' Variablen sind $A, B, C, D, E, F$.
\end{merkbox}

\subsection*{Schritt 3: Die Formel zusammenbauen}

F\"ur jeden Draht eine \"Aquivalenz. Das ergibt sechs Teilformeln,
die alle mit $\wedge$ verkn\"upft werden:

\begin{align*}
\Phi \;\equiv\; & \bigl(o_1 \leftrightarrow (A \wedge B)\bigr) \\
& \wedge\; \bigl(o_2 \leftrightarrow ((A \wedge \neg C) \vee (\neg A \wedge C))\bigr) \\
& \wedge\; \bigl(o_3 \leftrightarrow (B \vee C)\bigr) \\
& \wedge\; \bigl(D \leftrightarrow \neg o_1\bigr) \\
& \wedge\; \bigl(E \leftrightarrow (o_1 \wedge o_2 \wedge o_3)\bigr) \\
& \wedge\; \bigl(F \leftrightarrow (o_2 \wedge o_3)\bigr).
\end{align*}

\begin{analogie}
Jede Zeile der Formel ist wie ein \textbf{Vertragsparagraph}. "`Der
Draht $o_1$ f\"uhrt genau dann Strom, wenn sowohl $A$ als auch $B$
Strom f\"uhren."' Jeder Paragraph muss gleichzeitig gelten (das
$\wedge$ zwischen den Zeilen). Nur wenn alle Paragraphen erf\"ullt
sind, ist die Gesamtkonfiguration des Schaltkreises zul\"assig.
\end{analogie}

\subsection*{Schritt 4: Die drei Fragen beantworten}

Alle drei Fragen sind Fragen der Art "`Kann \ldots gelten?"' oder
"`Gilt immer \ldots?"'. Wir l\"osen sie mit dem Erf\"ullbarkeitstester
(siehe Teil~0.7).

\begin{center}
\renewcommand{\arraystretch}{1.4}
\begin{tabular}{c|l|l}
\toprule
& \textbf{Frage} & \textbf{Dem Tester geben} \\
\midrule
(a) & Kann $F = 1$ sein? & $\Phi \wedge F$ auf SAT testen \\
(b) & Kann $E = A$ sein? & $\Phi \wedge (E \leftrightarrow A)$ auf SAT testen \\
(c) & Gilt $F = 1 \Leftrightarrow D = 0$? & $\Phi \wedge \neg(F \leftrightarrow \neg D)$ auf UNSAT testen \\
\bottomrule
\end{tabular}
\end{center}

Statt einen Solver laufen zu lassen, k\"onnen wir es auch direkt mit
konkreten Belegungen nachpr\"ufen -- das ist \"ubersichtlicher:

\subsubsection*{(a) Kann $F = 1$ sein?}

\textbf{Methode.} Wir suchen eine Belegung $(A, B, C)$, f\"ur die
sich $F = 1$ ergibt.

\textbf{Versuch: $A = 1, B = 1, C = 0$.} Dann:
\begin{itemize}[nosep]
\item $o_1 = A \wedge B = 1 \wedge 1 = 1$.
\item $o_2 = (A \wedge \neg C) \vee (\neg A \wedge C) =
      (1 \wedge 1) \vee (0 \wedge 0) = 1 \vee 0 = 1$.
\item $o_3 = B \vee C = 1 \vee 0 = 1$.
\item $F = o_2 \wedge o_3 = 1 \wedge 1 = 1$. \checkmark
\end{itemize}

\textbf{Antwort:} \fbox{Ja, $F = 1$ ist m\"oglich.}

\subsubsection*{(b) Kann $E = A$ sein?}

\textbf{Methode.} Wir suchen eine Belegung, in der $E$ und $A$
denselben Wert haben.

\textbf{Trivial-Versuch: $A = 0, B = 0, C = 0$.} Dann:
\begin{itemize}[nosep]
\item $o_1 = 0 \wedge 0 = 0$, $o_2 = (0 \wedge 1) \vee (1 \wedge 0) = 0$,
      $o_3 = 0 \vee 0 = 0$.
\item $E = o_1 \wedge o_2 \wedge o_3 = 0 = A$. \checkmark
\end{itemize}

Das ist aber ein billiger Sieg. Spannender: gibt es auch eine
Belegung mit $A = 1$ und $E = 1$?

\textbf{Versuch: $A = 1, B = 1, C = 0$.} (Dieselbe wie bei (a).)
\begin{itemize}[nosep]
\item $o_1 = 1, o_2 = 1, o_3 = 1$.
\item $E = o_1 \wedge o_2 \wedge o_3 = 1 = A$. \checkmark
\end{itemize}

\textbf{Antwort:} \fbox{Ja, $E = A$ ist m\"oglich -- sogar auf zwei verschiedene Arten.}

\subsubsection*{(c) Gilt $F = 1 \Leftrightarrow D = 0$?}

Das ist eine \emph{immer}-Frage. Wir fragen: Stimmt es, dass in
\emph{jedem} zul\"assigen Zustand $F$ und $\neg D$ \"ubereinstimmen?

\textbf{Methode.} Wenn die \"Aquivalenz immer gilt, dann ist ihr
Gegenteil nie erf\"ullbar, also $\Phi \wedge \neg(F \leftrightarrow
\neg D)$ unerf\"ullbar. Findet man ein \textbf{Gegenbeispiel}, dann
gilt die \"Aquivalenz nicht.

\textbf{Suche nach Gegenbeispiel: $A = 1, B = 1, C = 1$.}
\begin{itemize}[nosep]
\item $o_1 = A \wedge B = 1 \wedge 1 = 1$. Also $D = \neg o_1 = 0$.
\item $o_2 = (A \wedge \neg C) \vee (\neg A \wedge C) =
      (1 \wedge 0) \vee (0 \wedge 1) = 0$.
\item $o_3 = B \vee C = 1 \vee 1 = 1$.
\item $F = o_2 \wedge o_3 = 0 \wedge 1 = 0$.
\end{itemize}
In diesem Zustand gilt $D = 0$, aber $F = 0$. Die rechte Seite
$D = 0$ ist wahr, die linke Seite $F = 1$ ist falsch. Damit ist die
Biimplikation $F = 1 \Leftrightarrow D = 0$ verletzt.

\textbf{Antwort:} \fbox{Nein, die \"Aquivalenz gilt NICHT generell.}
Gegenbeispiel: $(A, B, C) = (1, 1, 1)$, hier ist $D = 0$, aber
$F = 0$.

\begin{merkbox}[title={Intuition zum Gegenbeispiel}]
$D = 0$ hei{\ss}t, dass oben links beide Inputs $1$ sind ($A = B = 1$).
$F = 1$ dagegen braucht, dass beide Outputs $o_2$ und $o_3$ gleich $1$
sind -- und $o_2$ ist ein XOR-Gate \"uber $A$ und $C$, das wird $0$,
sobald $A = C$. Also: sobald $A = B = C$, widerspricht sich $D = 0$
und $F = 1$. Ein klares Gegenbeispiel.
\end{merkbox}

\newpage

% ============================================================
% AUFGABE 3
% ============================================================

\section{Aufgabe 3: Bishop Independence Problem}

\subsection*{Aufgabe in einfachen Worten}

Wir haben ein $n \times n$-Schachbrett. Gesucht sind alle
Konfigurationen, in denen L\"aufer (Bishops) so platziert sind, dass
\emph{keine zwei einander schlagen k\"onnen}. Zur Erinnerung: Der
L\"aufer zieht diagonal. Zwei L\"aufer schlagen sich also genau dann,
wenn sie auf derselben Diagonale stehen.

Der Auftrag: Schreibe ein Skript, das f\"ur gegebenes $n$ eine
aussagenlogische Formel produziert. Die Modelle der Formel sollen
genau die g\"ultigen Konfigurationen sein.

\subsection*{Schritt 1: Vergleich mit Rook und Queen}

In der Vorlesung wurden verwandte Probleme besprochen:

\begin{center}
\renewcommand{\arraystretch}{1.3}
\begin{tabular}{l|l}
\toprule
\textbf{Figur} & \textbf{Schl\"agt auf} \\
\midrule
Rook (Turm) & Zeile, Spalte \\
Bishop (L\"aufer) & Diagonalen (haupt + gegen) \\
Queen (Dame) & Zeile, Spalte, Diagonalen \\
\bottomrule
\end{tabular}
\end{center}

Die Encoding-Strategie ist bei allen drei gleich. Was sich \"andert,
sind die \textbf{Gruppen}, in denen der AMO-Constraint gilt:
\begin{itemize}[nosep]
\item Rook: AMO pro Zeile + AMO pro Spalte.
\item Bishop: AMO pro Hauptdiagonale + AMO pro Gegendiagonale.
\item Queen: AMO pro Zeile + Spalte + beide Diagonalen-Arten.
\end{itemize}

\begin{analogie}
Der L\"aufer ist ein \textbf{"`Diagonal-Turm"'}: Statt sich rechts,
links, oben oder unten zu bewegen, wandert er schr\"ag. Alles, was wir
f\"ur den Turm wussten (AMO pro Zeile/Spalte), \"ubertragen wir jetzt
auf die Diagonalen.
\end{analogie}

\subsection*{Schritt 2: Variablen und Diagonalen}

\begin{merkbox}[title={Die Variablen}]
F\"ur jedes Feld $(i, j)$ mit $1 \le i, j \le n$ definiere
\[
b_{i,j} \;:=\; \text{``Auf Feld $(i,j)$ steht ein L\"aufer''}.
\]
\end{merkbox}

Zur Erinnerung aus Teil~0.8: Jedes Feld $(i,j)$ liegt auf zwei
Diagonalen:
\begin{itemize}[nosep]
\item der Hauptdiagonale mit Index $i - j$,
\item der Gegendiagonale mit Index $i + j$.
\end{itemize}

\textbf{Beispielbrett $n = 4$ mit den Hauptdiagonalen eingezeichnet
(Zahl = Index $i-j$):}

\begin{center}
\renewcommand{\arraystretch}{1.3}
\begin{tabular}{|c|c|c|c|}
\hline
$0$ & $-1$ & $-2$ & $-3$ \\ \hline
$1$ & $0$ & $-1$ & $-2$ \\ \hline
$2$ & $1$ & $0$ & $-1$ \\ \hline
$3$ & $2$ & $1$ & $0$ \\ \hline
\end{tabular}
\end{center}

Felder mit dem gleichen Wert liegen auf derselben Hauptdiagonale --
z.\,B.\ die vier Felder mit $0$ bilden die Hauptdiagonale $(1,1),
(2,2), (3,3), (4,4)$.

\textbf{Dasselbe Brett mit Gegendiagonalen (Zahl = Index $i+j$):}

\begin{center}
\renewcommand{\arraystretch}{1.3}
\begin{tabular}{|c|c|c|c|}
\hline
$2$ & $3$ & $4$ & $5$ \\ \hline
$3$ & $4$ & $5$ & $6$ \\ \hline
$4$ & $5$ & $6$ & $7$ \\ \hline
$5$ & $6$ & $7$ & $8$ \\ \hline
\end{tabular}
\end{center}

Hier bilden z.\,B.\ die Felder mit $5$ die Gegendiagonale $(1,4),
(2,3), (3,2), (4,1)$.

\subsection*{Schritt 3: Die Constraints formulieren}

\begin{merkbox}[title={Bishop-Constraint}]
F\"ur jede Diagonale $D$ (egal ob Haupt oder Gegen) mit mindestens
zwei Feldern gilt: \textbf{h\"ochstens ein L\"aufer auf $D$}. In
paarweiser Enkodierung (AMO):
\[
\bigwedge_{\substack{(i_1,j_1), (i_2,j_2) \in D \\ (i_1,j_1) \neq (i_2,j_2)}}
\neg\bigl(b_{i_1,j_1} \wedge b_{i_2,j_2}\bigr).
\]
\end{merkbox}

Die Gesamtformel ist die Konjunktion aller dieser Klauseln, \"uber
alle Diagonalen hinweg.

\textbf{Warum keine Bedingung "`mindestens ein L\"aufer"'?} Die
Aufgabe fragt nur nach \emph{Unabh\"angigkeit}, nicht nach einer
Mindestanzahl. Das leere Brett ist eine triviale L\"osung (keine zwei
L\"aufer schlagen sich, wenn es gar keine gibt). Wenn wir zus\"atzlich
verlangen wollten, dass mindestens $k$ L\"aufer stehen, m\"ussten wir
daf\"ur extra Constraints hinzuf\"ugen.

\subsection*{Schritt 4: Das Python-Skript}

Das Skript geht in zwei Phasen vor: erst alle Diagonalen sammeln,
dann pro Diagonale die AMO-Klauseln erzeugen.

\begin{lstlisting}
def bishop_independence(n):
    """
    Erzeugt die aussagenlogische Formel (als Liste von Klauseln) fuer
    das Bishop Independence Problem auf einem n x n Schachbrett.

    Variablen: b_{i,j} (1 <= i,j <= n) steht fuer
    "Laeufer auf Feld (i,j)".
    Constraint: Fuer jede Diagonale hoechstens ein Laeufer
    (paarweise Klauseln, AMO-Enkodierung).
    """
    # Schritt 1: Alle Diagonalen sammeln.
    # Schluessel: (typ, index) -> Liste der Felder auf dieser Diagonale.
    diagonals = {}
    for i in range(1, n+1):
        for j in range(1, n+1):
            # Hauptdiagonale: i - j ist konstant.
            diagonals.setdefault(('haupt', i-j), []).append((i, j))
            # Gegendiagonale: i + j ist konstant.
            diagonals.setdefault(('gegen', i+j), []).append((i, j))

    # Schritt 2: Paarweise AMO-Klauseln erzeugen.
    clauses = []
    for feld_liste in diagonals.values():
        if len(feld_liste) < 2:
            continue  # einelementige Diagonalen brauchen keinen Constraint
        for p in range(len(feld_liste)):
            for q in range(p+1, len(feld_liste)):
                (i1, j1) = feld_liste[p]
                (i2, j2) = feld_liste[q]
                klausel = f"(NOT b_{{{i1},{j1}}} OR NOT b_{{{i2},{j2}}})"
                clauses.append(klausel)
    return clauses


def formel_ausgabe(clauses):
    """Verbindet alle Klauseln mit AND zur Gesamtformel."""
    return " AND ".join(clauses)


if __name__ == "__main__":
    n = 4
    klauseln = bishop_independence(n)
    print(f"Anzahl Klauseln fuer n={n}: {len(klauseln)}")
    for k in klauseln[:5]:
        print(k)
    print("...")
\end{lstlisting}

\subsection*{Schritt 5: Das Skript Zeile f\"ur Zeile}

\textbf{Das Dictionary \texttt{diagonals}.}
\begin{itemize}[nosep]
\item Wir benutzen ein Python-Dictionary (ein "`Nachschlagewerk"'), in
      dem der Schl\"ussel eine Diagonale identifiziert und der Wert
      eine Liste der Felder dieser Diagonale ist.
\item Der Schl\"ussel ist ein Tupel $(\text{typ}, \text{index})$,
      z.\,B.\ \texttt{('haupt', 0)} f\"ur die Hauptdiagonale mit
      $i - j = 0$.
\item \texttt{setdefault(key, [])} liefert die Liste zum Schl\"ussel
      zur\"uck -- und erzeugt sie, falls der Schl\"ussel noch nicht
      existiert, mit dem Vorgabewert \texttt{[]} (leere Liste). Danach
      h\"angen wir mit \texttt{.append((i,j))} das aktuelle Feld an.
\end{itemize}

\begin{analogie}
\texttt{diagonals.setdefault(key, []).append(...)} ist eine klassische
Python-Technik zum \textbf{Gruppieren}. Stell dir vor, du w\"urfst
Briefumschl\"age in Schubladen mit Etikett. F\"ur jeden Brief (Feld)
schaut man, ob die Schublade mit dem passenden Etikett
(\texttt{('haupt', 0)}) schon existiert. Wenn ja, rein damit. Wenn
nein, neue Schublade anlegen und reinlegen.
\end{analogie}

\textbf{Die doppelte Schleife \"uber $i$ und $j$.} Wir gehen jedes
Feld $(i,j)$ genau einmal durch. Pro Feld tragen wir es in zwei
Schubladen ein: einmal in die passende Hauptdiagonale, einmal in die
Gegendiagonale. Am Ende der Schleife hat jede Diagonale ihre Liste
von Feldern.

\textbf{Die Klauselerzeugung.} F\"ur jede Schublade (also jede
Diagonale):
\begin{itemize}[nosep]
\item Wenn die Schublade nur ein Feld enth\"alt (z.\,B.\ die Ecke
      $(1,4)$ ist auf der Hauptdiagonale $i-j = -3$ allein), gibt es
      nichts auszuschlie{\ss}en. Wir \"uberspringen sie mit
      \texttt{continue}.
\item Sonst generieren wir alle Paare $(p, q)$ mit $p < q$ (das ist
      eine Standardkonstruktion "`paarweise ohne Doppel"'). Jedes Paar
      gibt eine Klausel der Form
      \texttt{(NOT b\_\{i1,j1\} OR NOT b\_\{i2,j2\})}.
\end{itemize}

\textbf{Warum \texttt{NOT ... OR NOT ...}?} Die \"Ubersetzung: $\neg
(b_{i_1,j_1} \wedge b_{i_2,j_2})$ ist nach De Morgan dasselbe wie
$\neg b_{i_1,j_1} \vee \neg b_{i_2,j_2}$. Beide Schreibweisen sind
\"aquivalent, aber die zweite liegt bereits in CNF vor.

\subsection*{Schritt 6: Konkretes Beispiel $n = 4$}

F\"ur $n = 4$ gibt es jeweils $2n - 1 = 7$ Hauptdiagonalen und $7$
Gegendiagonalen. Ihre Gr\"o{\ss}en sind $1, 2, 3, 4, 3, 2, 1$ in beide
Richtungen (die mittelste Diagonale hat $n = 4$ Felder, dann wird sie
zu den Ecken hin kleiner).

\textbf{Anzahl Klauseln pro Diagonale:} $\binom{k}{2}$ f\"ur
Gr\"o{\ss}e $k$.
\begin{itemize}[nosep]
\item Gr\"o{\ss}e $1$: $\binom{1}{2} = 0$ Klauseln (wird \"ubersprungen).
\item Gr\"o{\ss}e $2$: $\binom{2}{2} = 1$ Klausel.
\item Gr\"o{\ss}e $3$: $\binom{3}{2} = 3$ Klauseln.
\item Gr\"o{\ss}e $4$: $\binom{4}{2} = 6$ Klauseln.
\end{itemize}

Summe pro Richtung: $0 + 1 + 3 + 6 + 3 + 1 + 0 = 14$. Beide Richtungen
zusammen: $14 + 14 = \mathbf{28}$ Klauseln.

\textbf{Beispiel-Klausel.} F\"ur die Hauptdiagonale $i - j = 0$,
zu der die Felder $(1,1), (2,2), (3,3), (4,4)$ geh\"oren, erzeugt das
Skript die $6$ Klauseln
\begin{align*}
& \neg(b_{1,1} \wedge b_{2,2}) \;\wedge\; \neg(b_{1,1} \wedge b_{3,3}) \;\wedge\; \neg(b_{1,1} \wedge b_{4,4}) \\
& \wedge\; \neg(b_{2,2} \wedge b_{3,3}) \;\wedge\; \neg(b_{2,2} \wedge b_{4,4}) \;\wedge\; \neg(b_{3,3} \wedge b_{4,4}).
\end{align*}
Lesart: "`Es stehen nicht gleichzeitig L\"aufer auf $(1,1)$ und
$(2,2)$"', "`nicht gleichzeitig auf $(1,1)$ und $(3,3)$"' usw.

\subsection*{Schritt 7: Was hei{\ss}t "`Modell der Formel"'?}

Ein Modell ist eine Belegung aller $b_{i,j}$ mit $0$ oder $1$. Eine
$1$ steht f\"ur "`hier steht ein L\"aufer"', eine $0$ f\"ur "`hier
steht kein L\"aufer"'.

\textbf{Triviales Modell:} alle $b_{i,j} = 0$. Dann ist jede Klausel
der Form $\neg(0 \wedge 0) = \neg 0 = 1$, also alle Klauseln wahr
-- das leere Brett ist erlaubt.

\textbf{Nicht-triviales Modell f\"ur $n = 4$:} Platziere L\"aufer auf
$(1,1), (1,3), (4,2), (4,4)$. Pr\"ufe:
\begin{itemize}[nosep]
\item $(1,1)$: Hauptdiagonale $i-j = 0$, Gegendiagonale $i+j = 2$.
\item $(1,3)$: Hauptdiagonale $i-j = -2$, Gegendiagonale $i+j = 4$.
\item $(4,2)$: Hauptdiagonale $i-j = 2$, Gegendiagonale $i+j = 6$.
\item $(4,4)$: Hauptdiagonale $i-j = 0$ -- \textbf{Kollision mit
      $(1,1)$!}
\end{itemize}
Diese Konfiguration ist also \emph{kein} Modell, weil $(1,1)$ und
$(4,4)$ auf derselben Hauptdiagonale liegen. Die Formel verbietet
das durch $\neg(b_{1,1} \wedge b_{4,4})$.

\textbf{Eine g\"ultige Konfiguration:} L\"aufer auf $(1,1), (1,4),
(4,1), (4,4)$.
\begin{itemize}[nosep]
\item Hauptdiagonalen: $1-1=0, 1-4=-3, 4-1=3, 4-4=0$.
      \textbf{Kollision: $(1,1)$ und $(4,4)$ auf $0$.}
\end{itemize}
Auch nicht erlaubt. L\"aufer auf den Ecken teilen sich n\"amlich die
beiden Hauptdiagonalen und Gegendiagonalen paarweise. Es ist
gar nicht so leicht, viele L\"aufer gleichzeitig unterzubringen!

\begin{merkbox}[title={Maximale L\"aufer-Anzahl}]
Auf einem $n \times n$-Brett kann man maximal $2n - 2$ unabh\"angige
L\"aufer platzieren (f\"ur $n \ge 2$). Unsere Formel hat aber auch
Modelle mit weniger L\"aufern, bis hin zum leeren Brett. Sie
beschreibt nur die \emph{Vertr\"aglichkeit}, nicht die Maximalit\"at.
\end{merkbox}

\begin{merkbox}[title={Pattern zum Mitnehmen}]
Das Python-Idiom
\begin{center}
\texttt{dict.setdefault(key, []).append(value)}
\end{center}
ist Standard zum \textbf{Gruppieren von Objekten nach Schl\"ussel}.
Es ersetzt umst\"andliches \texttt{if key not in dict: dict[key] = [];
dict[key].append(value)} durch einen einzigen Einzeiler.
\end{merkbox}

\newpage

% ============================================================
% ZUSAMMENFASSUNG
% ============================================================

\section*{Zusammenfassung und goldene Regeln}

\subsection*{Kurz\"ubersicht der Ergebnisse}

\begin{center}
\renewcommand{\arraystretch}{1.4}
\begin{tabular}{c|l|l}
\toprule
\textbf{Aufgabe} & \textbf{Inhalt} & \textbf{Wichtigstes Ergebnis} \\
\midrule
1 & Sudoku als Formel & $64$ Variablen, $400$ Klauseln, $5$ Constraint-Arten \\
2 & Schaltkreis als Formel & Hilfsvariablen $o_1, o_2, o_3$; 6 \"Aquivalenzen \\
2a & $F = 1$ m\"oglich? & Ja, z.\,B.\ $(A,B,C) = (1,1,0)$ \\
2b & $E = A$ m\"oglich? & Ja, z.\,B.\ $(A,B,C) = (1,1,0)$ oder $(0,0,0)$ \\
2c & $F = 1 \Leftrightarrow D = 0$? & Nein, Gegenbeispiel $(1,1,1)$ \\
3 & Bishop Independence & Skript erzeugt AMO-Klauseln pro Diagonale \\
3 ($n=4$) & Anzahl Klauseln & $28$ (je $14$ Haupt-/Gegendiagonalen) \\
\bottomrule
\end{tabular}
\end{center}

\subsection*{Goldene Regeln}

\begin{merkbox}[title={Goldene Regeln zum Mitnehmen}]
\begin{enumerate}
\item \textbf{Encoding-Rezept.} Eine Bool'sche Variable pro
      Entscheidung. Constraints als Konjunktion von
      \textbf{AMO-}, \textbf{ALO-} und \textbf{EO}-Teilformeln.
      Gesamtformel = Konjunktion aller Regeln.

\item \textbf{AMO paarweise.} F\"ur "`h\"ochstens eine von $k$
      Variablen wahr"' schreibt man
      $\bigwedge_{i<j} \neg(x_i \wedge x_j)$. Anzahl: $\binom{k}{2}$
      Klauseln. Funktioniert immer, wenn $k$ nicht zu gro{\ss} wird.

\item \textbf{EO = ALO $\wedge$ AMO.} "`Genau eine"' zerf\"allt in
      "`mindestens eine"' (eine gro{\ss}e Disjunktion) und "`h\"ochstens
      eine"' (viele Paar-Ausschl\"usse).

\item \textbf{Schaltkreise.} Pro Gate-Output eine \"Aquivalenz, alle
      mit $\wedge$ verbunden. F\"ur interne Dr\"ahte Hilfsvariablen
      einf\"uhren. Summen-Bedingungen wie $[>k], [=k]$ mit
      Fallunterscheidung in Bool'sche Ausdr\"ucke \"ubersetzen.

\item \textbf{Fragen an den Erf\"ullbarkeitstester.}
      \begin{itemize}[nosep]
      \item "`Kann $X$ passieren?"' $\Rightarrow$ SAT-Test auf
            $\Phi \wedge X$.
      \item "`Gilt immer $Y$?"' $\Rightarrow$ UNSAT-Test auf
            $\Phi \wedge \neg Y$.
      \end{itemize}

\item \textbf{Schach-Diagonalen.} Hauptdiagonale: $i - j$ konstant
      (links oben $\to$ rechts unten). Gegendiagonale: $i + j$
      konstant (rechts oben $\to$ links unten). Je $2n - 1$ St\"uck
      auf einem $n \times n$-Brett.

\item \textbf{Sanity-Check.} Anzahl der Klauseln und Variablen
      mitz\"ahlen. Falls die Zahlen nicht passen, irgendwo fehlt eine
      Bedingung oder es gibt Doppelz\"ahlungen.
\end{enumerate}
\end{merkbox}

\begin{analogie}[title={Leitmotiv}]
Aussagenlogik ist eine \textbf{universelle Puzzle-Sprache}. Jedes
Problem aus Kombinatorik, Schach oder Schaltkreisbau l\"asst sich in
diese Sprache \"ubersetzen. Einmal \"ubersetzt, steht ein ganzer
Baukasten an Werkzeugen bereit: SAT-Solver, \"Aquivalenzumformungen,
Normalformen. Das ist der Grund, warum Logik so viel mehr ist als
nur eine mathematische Spielerei -- sie ist die Grundlage, auf der
Computer Probleme "`verstehen"' und l\"osen k\"onnen.
\end{analogie}

\bigskip
\bigskip

Damit sind alle drei Aufgaben des Blattes durchgearbeitet. Wer die
drei Schritte \emph{Variablen einf\"uhren}, \emph{Regeln als
Teilformeln}, \emph{alles mit $\wedge$ verbinden} verinnerlicht hat,
kann damit weit \"uber Sudoku, Schach und Schaltkreise hinaus gehen.
Viel Spa{\ss} beim Weiter-Encodieren!

\end{document}
